\section{奇异值分解}

\subsection{奇异值分解}

\begin{BoxDefinition}[奇异值与奇异向量]
    对于$\vb{A}\in\R^{m\times n}$，若存在$\vb{u}\in\R^m,\vb{v}\in\R^n$且$\vb{u},\vb{v}\neq\vb{0}$以及$\sigma\in\R_{+}$，使
    \begin{Equation}[奇异值与奇异向量]
        \vb{A}\vb{v}=\sigma\vb{u}\qquad
        \vb{A}^T\vb{u}=\sigma\vb{v}
    \end{Equation}
    则称$\sigma$是奇异值（Singular Value），而称$\vb{u},\vb{v}$为左右奇异向量（Singular Vector）。
\end{BoxDefinition}

\begin{BoxDefinition}[奇异值分解]
    对于矩阵$\vb{A}\in\R^{m\times n}$，定义其奇异值分解（Singular Value Decomposition, SVD）
    \begin{Equation}[奇异值分解]
        \vb{A}=\vb{U}\vb*{\Sigma}\vb{V}^T
    \end{Equation}
    其中，$\vb{U}\in\R^{m\times m},\vb{V}\in\R^{n\times n}$是正交矩阵，$\vb*{\Sigma}\in\R^{m\times n}$是对角矩阵
    \begin{Equation}[奇异值分解的对角矩阵]
        \vb*{\Sigma}_{ij}=\begin{cases}
            \sigma_i,&i=j\\
            0,&i\neq j\\
        \end{cases}
    \end{Equation}
    且所有元素按顺序排列并均非负$\sigma_1\geq\sigma_2\geq\cdots\geq\sigma_p\geq 0$，其中$p=\min\qty{m,n}$。
\end{BoxDefinition}

\begin{BoxTheorem}[奇异值分解的存在性]
    $\vb{A}\in\R^{m\times n}$总是存在奇异值分解。
\end{BoxTheorem}

奇异值分解$\vb{A}=\vb{U}\vb*{\Sigma}\vb{V}^T$和特征值分解$\vb{A}=\vb{V}\vb*{\Lambda}\vb{V}^{-1}$在形式上是很相似的，从某种意义上可以视为特征值分解的扩展。首先，仅有可对角化的方阵$\R^{n\times n}$存在特征值分解，但任意的一般矩阵$\R^{m\times n}$都总是存在奇异值分解，适用性提高了很多！其次，特征值和特征向量的概念是来自$\vb{A}\vb{v}=\lambda\vb{v}$的关系，奇异值和奇异向量则将其裂分为相互耦合$\vb{A}\vb{v}=\sigma\vb{u}$和$\vb{A}^T\vb{u}=\sigma\vb{v}$的左右两组关系，即$\vb{A}$和$\vb{A}^T$可以在右奇异向量$\vb{v}$和左奇异向量$\vb{u}$间来回变换，且效果是相同的非负常数倍$\sigma>0$。容易证明，$\vb{U}$、$\vb{V}$、$\vb*{\Sigma}$提供了$i=1,\cdots,p$组奇异向量和奇异值$\vb{u}_i、\vb{v}_i、\sigma_i$。

对$\vb{A}=\vb{U}\vb*{\Sigma}\vb{V}^T$两边右乘$\vb{V}$，利用$\vb{V}$的正交性
\begin{Equation}[奇异讨论1]
    \vb{A}\vb{V}=\vb{U}\vb*{\Sigma}\qquad \vb{A}\vb{v}_i=\sigma_i\vb{u}_i,\ i=1,\cdots,p
\end{Equation}

对$\vb{A}^T=\vb{V}\vb*{\Sigma}^T\vb{U}^T$两边右乘$\vb{U}$，利用$\vb{U}$的正交性
\begin{Equation}[奇异讨论2]
    \vb{A}^T\vb{U}=\vb{V}\vb*{\Sigma}^T\qquad \vb{A}^T\vb{u}_i=\sigma_i\vb{v}_i,\ i=1,\cdots,p
\end{Equation}

上述我们讨论的都是Full SVD，实际上可以写出更为紧凑的Compact SVD，这种压缩主要来自两个方面：首先，对于$m>n$和$m<n$的情况，由于$p=\min\qty{m,n}$，因此$\vb{U}$和$\vb{V}$在两种情况下都会有一部分不影响结果但为了凑正交方阵的列向量。其次，奇异值只规定非负但可以等于零，设有$r$个非负奇异值，即有$\sigma_1\geq\sigma_2\geq\cdots\geq\sigma_r>0$而$\sigma_{r+1}=\cdots=\sigma_p=0$。

\begin{BoxEquation}[奇异值分解的分块形式]
    设$\vb{A}\in\R^{m\times n}$有$r$个非零奇异值，则奇异值分解$\vb{A}=\vb{U}\vb*{\Sigma}\vb{V}^T$可以写为以下分块形式
    \begin{Equation}[奇异值分解的分块形式]
        \vb{A}=
        \mqty[\vb{U}_1&\vb{U}_2]
        \mqty[\tilde{\vb*{\Sigma}}&\vb{0}\\ \vb{0}&\vb{0}]
        \mqty[\vb{V}_1^T\\ \vb{V}_2^T]
    \end{Equation}
    或进一步写为
    \begin{Equation}[奇异值分解的压缩形式]
        \vb{A}=\vb{U}_1\vb*{\Sigma}\vb{V}^T
    \end{Equation}
    $\tilde{\vb*{\Sigma}}\in\R^{r\times r}$可以表示为
    \begin{Equation}[奇异值分解的分块形式1]
        \tilde{\vb*{\Sigma}}=\diag(\sigma_1,\cdots,\sigma_r)
    \end{Equation}
    $\vb{U}_1\in\R^{m\times r},\vb{U}_2\in\R^{m\times(m-r)}$可以表示为
    \begin{Equation}[奇异值分解的分块形式2]
        \vb{U}_1=\qty[\vb{u}_1,\cdots,\vb{u}_r]\qquad
        \vb{U}_2=\qty[\vb{u}_{r+1},\cdots,\vb{u}_{m}]
    \end{Equation}
    $\vb{V}_1\in\R^{n\times r},\vb{V}_2\in\R^{n\times(n-r)}$可以表示为
    \begin{Equation}[奇异值分解的分块形式3]
        \vb{V}_1=\qty[\vb{v}_1,\cdots,\vb{v}_r]\qquad
        \vb{V}_2=\qty[\vb{v}_{r+1},\cdots,\vb{v}_{n}]
    \end{Equation}
\end{BoxEquation}

\subsection{奇异分解和特征分解的关系}

奇异分解和特征分解的形式是十分相似的，有趣的是，对于半正定矩阵，它具有相同的奇异分解和特征分解！\cref{thm:厄米矩阵的特征分解}告诉我们对称矩阵总是存在特征分解，特征值为实数，特征向量相互正交，\cref{thm:半正定矩阵的构造}告诉我们半正定矩阵总是具有$\vb{A}\vb{A}^T$或$\vb{A}^T\vb{A}$的形式，因此有以下两个性质。

\begin{BoxProperty}[奇异分解和特征分解1]
    设$\vb{A}\in\R^{m\times n}$的奇异分解为$\vb{A}=\vb{U}\vb*{\Sigma}\vb{V}^T$，则$\vb{A}\vb{A}^T$具有相同的特征分解和奇异分解
    \begin{Gather}[奇异分解和特征分解1]
        \vb{A}\vb{A}^T=\vb{U}\vb{D}_1\vb{U}^T\in\R^{m\times m}\id{a}\\
        \vb{D}_1=\vb*{\Sigma}\vb*{\Sigma}^T=\diag(\sigma_1^2,\cdots,\sigma_p^2,0,\cdots,0)\in\R^{m\times m}\id{b}
    \end{Gather}
\end{BoxProperty}

\begin{BoxProperty}[奇异分解和特征分解2]
    设$\vb{A}\in\R^{m\times n}$的奇异分解为$\vb{A}=\vb{U}\vb*{\Sigma}\vb{V}^T$，则$\vb{A}^T\vb{A}$具有相同的特征分解和奇异分解
    \begin{Gather}[奇异分解和特征分解1]
        \vb{A}^T\vb{A}=\vb{V}\vb{D}_2\vb{V}^T\in\R^{n\times n}\id{a}\\
        \vb{D}_2=\vb*{\Sigma}^T\vb*{\Sigma}=\diag(\sigma_1^2,\cdots,\sigma_p^2,0,\cdots,0)\in\R^{n\times n}\id{b}
    \end{Gather}
\end{BoxProperty}

\subsubsection{奇异值分解的性质}



